\documentclass[a4paper,10pt,BCOR10mm,oneside,headsepline]{scrartcl}
\usepackage[portuges]{babel}
\usepackage[utf8]{inputenc}
\usepackage{amsmath}
\usepackage{wasysym}% provides \ocircle and \Box
\usepackage{enumitem}% easy control of topsep and leftmargin for lists
\usepackage{color}% used for background color
\usepackage{forloop}% used for \Qrating and \Qlines
\usepackage{ifthen}% used for \Qitem and \QItem
\usepackage{typearea}
\usepackage{graphicx}
\usepackage{booktabs}
\usepackage{listings}
\usepackage{color}
\usepackage{algorithm}
\usepackage{algorithmic}
\usepackage{amssymb}
\usepackage{minted}
\usepackage{mdframed}
\usepackage{wasysym}
\usepackage{xcolor}
\usepackage{fontspec}
%\usepackage{tgtermes}

\areaset{17cm}{26cm}
\setlength{\topmargin}{-2cm}

\definecolor{codegreen}{rgb}{0,0.6,0}
\definecolor{codegray}{rgb}{0.5,0.5,0.5}
\definecolor{codepurple}{rgb}{0.58,0,0.82}
\definecolor{backcolour}{rgb}{0.95,0.95,0.92}
 
\lstdefinestyle{mystyle}{
    backgroundcolor=\color{backcolour},   
    commentstyle=\color{codegreen},
    keywordstyle=\color{magenta},
    numberstyle=\tiny\color{codegray},
    stringstyle=\color{codepurple},
    basicstyle=\footnotesize,
    breakatwhitespace=false,         
    breaklines=true,                 
    captionpos=b,                    
    keepspaces=true,                 
    numbers=left,                    
    numbersep=5pt,                  
    showspaces=false,                
    showstringspaces=false,
    showtabs=false,                  
    tabsize=2
}
 
\lstset{style=mystyle}

\mdfsetup{
    skipabove=12pt,
    skipbelow=12pt,
    innertopmargin=8pt,
    innerbottommargin=8pt,
    innerleftmargin=10pt,
    innerrightmargin=10pt,
    splittopskip=12pt,
    splitbottomskip=12pt,
}

\algsetup{linenodelimiter=.,linenosize=\tiny}

\renewcommand{\algorithmicrequire}{\textbf{Entrada:}}
\renewcommand{\algorithmicensure}{\textbf{Saída:}}
\floatname{algorithm}{Algoritmo}

\setmainfont{TeX Gyre Termes}
\setmonofont{APL 385 Unicode}[Scale=0.9]

%% \Qq = Questionaire question. Oh, this is just too simple. It helps
%% making it easy to globally change the appearance of questions.
\newcommand{\Qq}[1]{\textbf{#1}}
%\newcommand{\Qq}[1]{{#1}}

%% \QO = Circle or box to be ticked. Used both by direct call and by
%% \Qrating and \Qlist.
\newcommand{\QO}{$\Box$}% or: $\ocircle$

%% \Qrating = Automatically create a rating scale with NUM steps, like
%% this: 0--0--0--0--0.
\newcounter{qr}
\newcommand{\Qrating}[1]{\QO\forloop{qr}{1}{\value{qr} < #1}{---\QO}}

%% \Qline = Again, this is very simple. It helps setting the line
%% thickness globally. Used both by direct call and by \Qlines.
\newcommand{\Qline}[1]{\noindent\rule{#1}{0.6pt}}

%% \Qlines = Insert NUM lines with width=\linewith. You can change the
%% \vskip value to adjust the spacing.
\newcounter{ql}
\newcommand{\Qlines}[1]{\forloop{ql}{0}{\value{ql}<#1}{\vskip0em\Qline{\linewidth}}}

%% \Qlist = This is an environment very similar to itemize but with
%% \QO in front of each list item. Useful for classical multiple
%% choice. Change leftmargin and topsep accourding to your taste.
\newenvironment{Qlist}{%
\renewcommand{\labelitemi}{\QO}
\begin{itemize}[leftmargin=1.5em,topsep=-.5em]
}{%
\end{itemize}
}

%% \Qtab = A "tabulator simulation". The first argument is the
%% distance from the left margin. The second argument is content which
%% is indented within the current row.
\newlength{\qt}
\newcommand{\Qtab}[2]{
\setlength{\qt}{\linewidth}
\addtolength{\qt}{-#1}
\hfill\parbox[t]{\qt}{\raggedright #2}
}

%% \Qitem = Item with automatic numbering. The first optional argument
%% can be used to create sub-items like 2a, 2b, 2c, ... The item
%% number is increased if the first argument is omitted or equals 'a'.
%% You will have to adjust this if you prefer a different numbering
%% scheme. Adjust topsep and leftmargin as needed.
\newcounter{itemnummer}
\newcommand{\Qitem}[2][]{% #1 optional, #2 notwendig
\ifthenelse{\equal{#1}{}}{\stepcounter{itemnummer}}{}
\ifthenelse{\equal{#1}{a}}{\stepcounter{itemnummer}}{}
\begin{enumerate}[topsep=2pt,leftmargin=2.8em]
\item[\textbf{\arabic{itemnummer}#1.}] #2
\end{enumerate}
}

%% \QItem = Like \Qitem but with alternating background color. This
%% might be error prone as I hard-coded some lengths (-5.25pt and
%% -3pt)! I do not yet understand why I need them.
\definecolor{bgodd}{rgb}{0.8,0.8,0.8}
\definecolor{bgeven}{rgb}{0.9,0.9,0.9}
\newcounter{itemoddeven}
\newlength{\gb}
\newcommand{\QItem}[2][]{% #1 optional, #2 notwendig
\setlength{\gb}{\linewidth}
\addtolength{\gb}{-5.25pt}
\ifthenelse{\equal{\value{itemoddeven}}{0}}{%
\noindent\colorbox{bgeven}{\hskip-3pt\begin{minipage}{\gb}\Qitem[#1]{#2}\end{minipage}}%
\stepcounter{itemoddeven}%
}{%
\noindent\colorbox{bgodd}{\hskip-3pt\begin{minipage}{\gb}\Qitem[#1]{#2}\end{minipage}}%
\setcounter{itemoddeven}{0}%
}
}

\newcommand{\rawcode}[1]{\texttt{#1}}
\newcommand{\code}[2]{\mintinline{#1}{#2}}
\newcommand{\file}[1]{\texttt{\underline{#1}}}

\renewcommand{\theFancyVerbLine}{\tiny \texttt{\textcolor{gray}{\arabic{FancyVerbLine}}}}

\newcommand{\inputcode}[2]{\begin{mdframed}\inputminted[fontsize=\footnotesize,linenos,numbersep=4pt]{#1}{#2}\end{mdframed}}

\newenvironment{solution}{%
    \begin{mdframed}%
        \noindent\textbf{Solu\c{c}\~ao:}\par\medskip
}{%
    \end{mdframed}%
}

\newcommand{\inputsnippet}[4]{\begin{mdframed}[shadow=true,shadowsize=3pt,shadowcolor=gray]\inputminted[fontsize=\footnotesize,linenos,firstline=#2,lastline=#3,numbersep=2pt]{#1}{#4}\end{mdframed}}

\newcommand{\inputsyntax}[2]{\inputminted[fontsize=\footnotesize]{#1}{#2}}



\begin{document}

\begin{center}
	\includegraphics[scale=0.45]{fga}
	% fga.ps: 1087x103 pixel, 72dpi, 38.35x3.63 cm, bb=0 0 1087 103
\end{center}


\vspace{0.1in}

\begin{center}
    \textbf{\huge Paradigmas de Programação}

    \vspace{0.1in}
    \textbf{\Large Máquinas de Turing: Lista 1}
\end{center}\vskip1em

\renewcommand{\Qq}[1]{{#1}}

\section*{Parte A - Enumerabilidade}

\Qitem{ \Qq{Defina uma função $f: \mathbb{N}\to A$, total ou parcial, que enumere os elementos do
        conjunto $A$ definido abaixo.}

    \vspace{0.1in}
    \begin{Qlist}
    \item[(a)] O conjunto dos múltiplos positivos de três, isto é,
            \[
                A = \{ 0, 3, 6, 9, \ldots \}
            \]

    \begin{solution}
    $
        f(n) = 3(n-1)
    $
    \end{solution}

    \item[(b)] O conjunto dos números primos, isto é,
            \[
                A = \{ 2, 3, 5, 7, 11, 13, 17, 19, 23, \ldots \}
            \]

    \begin{solution}
    \begin{itemize}
        \item \textbf{Total:} $f(n) = p_n$, sendo $p_n$ o $n$-ésimo número primo, com $p_1=2$, $p_2=3$, $p_3=5$, $\ldots$
        \item \textbf{Parcial:}
        \[
            f(n) = \left\{ \begin{array}{ll}
                        n, & \mbox{se}\ n\ \mbox{é primo} \\
                        \mbox{indefinida}, & \mbox{caso contrário}
                    \end{array}\right.
        \]
    \end{itemize}
    \end{solution}

    \item[(c)] O conjunto das notas musicais, isto é,
            \[
                A = \{ \mbox{dó, ré, mi, fá, sol, lá, si} \}
            \]

    \begin{solution}
    \[
        f(n) = \left\{ \begin{array}{ll}
                    \mbox{dó}, & \mbox{se}\ n = 1 \\
                    \mbox{ré}, & \mbox{se}\ n = 2 \\
                    \mbox{mi}, & \mbox{se}\ n = 3 \\
                    \mbox{fá}, & \mbox{se}\ n = 4 \\
                    \mbox{sol}, & \mbox{se}\ n = 5 \\
                    \mbox{lá}, & \mbox{se}\ n = 6 \\
                    \mbox{si}, & \mbox{se}\ n = 7 \\
                    \mbox{indefinida}, & \mbox{se}\ n > 7
                \end{array}\right.
    \]
    \end{solution}
    \end{Qlist}
}

\Qitem{ \Qq{Compute os valores abaixo, onde $G$ é a função que enumera os pares de inteiros 
    positivos segundo o zigue-zague de Cantor.}

    \begin{solution}
        $G(1)=(1,1)$, $G(2)=(1,2)$, $G(3)=(2,1)$, $G(4)=(1,3)$, $\ldots$:
        os pares são listados em ordem crescente da soma $s$ de suas coordenadas e, para a mesma
        soma, em ordem crescente da primeira coordenada. A diagonal de soma $s$ tem $s-1$ pares:
        $(1,s-1),(2,s-2),\ldots,(s-1,1)$. Logo, o número de pares com soma até $s$ é
        $1+2+\cdots+(s-1) = \frac{(s-1)s}{2}$:
        \begin{center}
        \begin{tabular}{c|cccccccccccccc}
        $s$ & 2 & 3 & 4 & 5 & 6 & 7 & 8 & 9 & 10 & 11 & 12 & 13 & 14 & 15 \\ \midrule
        pares até $s$ & 1 & 3 & 6 & 10 & 15 & 21 & 28 & 36 & 45 & 55 & 66 & 78 & 91 & 105
        \end{tabular}
        \end{center}
        Para achar $G(n)$, encontra-se a menor soma $s$ com $\frac{(s-1)s}{2}\geq n$. Então
        $G(n)=(k,\,s-k)$, onde $k = n - \frac{(s-2)(s-1)}{2}$ é a posição de $n$ dentro da
        diagonal.
    \end{solution}

    \vspace{0.1in}
    \begin{Qlist}
        \item[(a)] $G(25)$

        \begin{solution}
            Até $s=7$ há $21$ pares; $s=8$ cobre as posições $22$ a $28$. Então
            $k = 25 - 21 = 4$ e
            \[
                G(25) = (4,\,8-4) = (4,\,4).
            \]
        \end{solution}

        \item[(b)] $G(50)$

        \begin{solution}
            Até $s=10$ há $45$ pares; $s=11$ cobre as posições $46$ a $55$. Então
            $k = 50 - 45 = 5$ e
            \[
                G(50) = (5,\,11-5) = (5,\,6).
            \]
        \end{solution}

        \item[(c)] $G(100)$

        \begin{solution}
            Até $s=14$ há $91$ pares; $s=15$ cobre as posições $92$ a $105$. Então
            $k = 100 - 91 = 9$ e
            \[
                G(100) = (9,\,15-9) = (9,\,6).
            \]
        \end{solution}
    \end{Qlist}
}

\Qitem{ \Qq{Em relação à enumeração das sequências finitas de inteiros positivos, determine o 
    natural que codifica a sequência indicada.}

    \begin{solution}
    $s=\{s_1,\ldots,s_N\}$ é codificada por
    $f(s) = 2^{s_1}3^{s_2}5^{s_3}\cdots p_N^{s_N}$, onde $p_N$ é o $N$-ésimo primo.
    \end{solution}

    \vspace{0.1in}
    \begin{Qlist}
        \item[(a)] $s_1 = \{ 1, 2, 1, 2, 1 \}$

            \begin{solution}
               $f(s_1) = 2^1\cdot 3^2\cdot 5^1\cdot 7^2\cdot 11^1
                = 2\cdot 9\cdot 5\cdot 49\cdot 11 = 48\,510.$ 
            \end{solution}

        \item[(b)] $s_2 = \{ 1, 2, 3, 4, 5 \}$

            \begin{solution}
               $f(s_2) = 2^1\cdot 3^2\cdot 5^3\cdot 7^4\cdot 11^5
                = 2\cdot 9\cdot 125\cdot 2\,401\cdot 161\,051 = 870\,037\,764\,750.$
            \end{solution}
        \item[(c)] $s_3 = \{ 1, 1, 1, 1, 1 \}$

            \begin{solution}
               $f(s_3) = 2^1\cdot 3^1\cdot 5^1\cdot 7^1\cdot 11^1 = 2\,310.$
            \end{solution}
    \end{Qlist}
}

\Qitem{ \Qq{Ainda em relação à enumeração das sequências finitas de inteiros positivos, 
        determine qual seria a $n$-ésima sequência, segundo esta enumeração.}

        \begin{solution}
            $G=f^{-1}$ é parcial: um natural $m$ decodifica uma sequência se
            $m=1$ ou se, sendo $m=2^{a_1}3^{a_2}\cdots p_k^{a_k}$, tem-se $a_i>0$
            para todo $i$.
            \begin{center}
            \begin{tabular}{c|cccccccccccc}
            posição & 1 & 2 & 3 & 4 & 5 & 6 & 7 & 8 & 9 & 10 & 11 & 12 \\ \midrule
            código $m$ & 1 & 2 & 4 & 6 & 8 & 12 & 16 & 18 & 24 & 30 & 32 & 36 \\
            sequência & $e$ & $\{1\}$ & $\{2\}$ & $\{1,1\}$ & $\{3\}$ & $\{2,1\}$ & $\{4\}$ &
                $\{1,2\}$ & $\{3,1\}$ & $\{1,1,1\}$ & $\{5\}$ & $\{2,2\}$
            \end{tabular}
            \end{center}
        \end{solution}

    \vspace{0.1in}
    \begin{Qlist}
        \item[(a)] $n = 8$

            \begin{solution}
                \item[(a)] $n=8$: $\{1, 2\}$, pois $18 = 2^1\cdot 3^2$.
            \end{solution}

        \item[(b)] $n = 10$

            \begin{solution}
                \item[(b)] $n=10$: $\{1, 1, 1\}$, pois $30 = 2^1\cdot 3^1\cdot 5^1$.
            \end{solution}
        \item[(c)] $n = 12$

            \begin{solution}
                \item[(c)] $n=12$: $\{2, 2\}$, pois $36 = 2^2\cdot 3^2$.
            \end{solution}
    \end{Qlist}
}

\Qitem{ \Qq{Dada uma lista $L = S_1, S_2, S_3, \ldots$ de conjuntos de inteiros positivos, 
    determine o conjunto $\Delta(L)$, de acordo com a diagonalização de Cantor.}

    \begin{solution}
        Dada $L = S_1, S_2, S_3, \ldots$, o conjunto diagonal é definido por
        \[
            \Delta(L) = \{ n\in\mathbb{Z}^+\ |\ n\not\in S_n \}.
        \]
    \end{solution}

    \vspace{0.1in}
    \begin{Qlist}
        \item[(a)] $S_i$ é o conjunto dos $i$ primeiros naturais:
            \[
                S_1 = \{ 1 \}, S_2 = \{ 1, 2 \}, S_3 = \{ 1, 2, 3 \}, \ldots
            \]

        \begin{solution}
            Como $S_n = \{1, \ldots, n\}$, tem-se $n \in S_n$ para todo $n \in \mathbb{Z}^+$. Logo,
            \[
                \Delta(L) = \emptyset.
            \]
        \end{solution}

        \item[(b)] $S_i$ é o conjunto dos divisores próprios de $i$, isto é, divisores positivos e 
            menores do que $i$:
            \[
                S_1 = \emptyset, S_2 = \{ 1 \}, S_3 = \{ 1 \}, S_4 = \{ 1, 2 \}, \ldots
            \]

        \begin{solution}
            Todos os elementos de $S_n$ são menores do que $n$, logo $n \not\in S_n$ para todo
            $n \in \mathbb{Z}^+$ (inclusive $n = 1$, pois $S_1 = \emptyset$). Assim,
            \[
                \Delta(L) = \mathbb{Z}^+ = \{1, 2, 3, \ldots\}.
            \]
        \end{solution}

        \item[(c)] $S_i$ é o conjunto dos $i$ primeiros números pares, isto é,
            \[
                S_1 = \{ 2 \}, S_2 = \{ 2, 4 \}, S_3 = \{ 2, 4 ,6 \}, \ldots 
            \]

        \begin{solution}
            Como $S_n = \{2, 4, \ldots, 2n\}$, tem-se $n \in S_n$ se $n$ é par ($n \leq 2n$), e
            $n \not\in S_n$ se $n$ é ímpar. Logo,
            \[
                \Delta(L) = \{1, 3, 5, 7, \ldots\}.
            \]
        \end{solution}
    \end{Qlist}
}



\end{document}
